%! Author = chtompki
%! Date = 1/14/26
\documentclass[12pt]{amsart}
% Default margins are too wide all the way around. I reset them here
\setlength{\hoffset}{-.75in}
\setlength{\textwidth}{6.5in}
\setlength{\voffset}{-.5in}
\setlength{\textheight}{9.0in}
\usepackage{amsmath}
\usepackage{leftindex}
\usepackage{amssymb}
\usepackage{amsthm}
\usepackage{enumitem}
\usepackage{setspace}

\newenvironment{solution}
{\begin{proof}[Solution]}
{\end{proof}}

\title{Combinatorics\\Assignment 2}
\author{Rob Tompkins}
\date{\today}

\begin{document}
    \maketitle
    \date{\today}
    \begin{onehalfspace}
        % Problem 6
        \noindent\textbf{6.} Prove that $$P(n,k) = k!\binom{n}{k}.$$
        \begin{proof}
            First let's consider how we choose our $k$-permutations of $[n]$.
            We begin by choosing the first element with $n$ choices, the second with $n-1$ choices, then all the way down to the $k^{\text{th}}$ element with $n-k+1$ choices, so we see that
            \begin{displaymath}
                P(n,k) = n(n-1)(n-2)\cdots(n-k+1) = \frac{n!}{(n-k)!}.
            \end{displaymath}
            On the right-hand side, we are choosing $k$-elements from $[n]$ in an unsorted fashion, but then
            we need to sort them in $k$ ways which is $k!$, and we see that
            \begin{displaymath}
                k!\binom{n}{k} = k!\left(\frac{n!}{k!(n-k)!}\right) = \frac{n!}{(n-k)!}.
            \end{displaymath}
        \end{proof}

        % Problem 10
        \textbf{10.} Prove that:
        \begin{displaymath}\sum_{k=0}^n r^{n-k}\binom{n}{k} = (r + 1)^n
        \end{displaymath}
        \begin{proof}
            We'll go ahead and consider the general case here for $(x+y)^n$.
            We want to proove:
            \begin{displaymath}
                (x+y)^n = \sum_{k=0}^n\binom{n}{k}x^{n-k}y^k
            \end{displaymath}
            which if we multiply out, we get $2^n$ terms each of which can be arranged as $x^{n-k}y^k$ for $k\in[n]$.
            We obtain the coefficient for the $x^{n-k}y^k$ term by thinking about each of our $2^n$ terms written out
            without using exponents.
            That is, that each term is a string of length $n$ of individual $x$'s and $y$'s (which are allowed to repeat).
            Now our goal is to count up the terms number of $x$'s is a specific value and the number of $y$'s is a specific value.
            Let the number of $y$'s appearing be $k$ for convenience.
            Then we are chosing the number of $k$-subsets from $[n]$, which we know to be $\binom{n}{k}$.
            So we have
            \begin{displaymath}
                (x + y)^n = \sum_{k=0}^n\binom{n}{k}x^{n-k}y^k.
            \end{displaymath}
            For the case in the problem we simply set $x=r$ and $y=1$, and we are done.

            \paragraph{\emph{Proof by induction of the Binomial Theorem.}}

            When \(n=0\), both sides equal \(1\), since \(x^{0}=1\) and
            \[
                \binom{0}{0}=1.
            \]

            Now suppose the equality holds for a given \(n\); we will prove it for \(n+1\).
            For \(j,k \ge 0\), let \([f(x,y)]_{j,k}\) denote the coefficient of \(x^{j}y^{k}\)
            in the polynomial \(f(x,y)\).
            By the inductive hypothesis, \((x+y)^n\) is a polynomial in \(x\) and \(y\) such that
            \[
                [(x+y)^n]_{j,k} =
                \begin{cases}
                    \binom{n}{k}, & \text{if } j+k=n,\\
                    0, & \text{otherwise}.
                \end{cases}
            \]

            The identity
            \[
                (x+y)^{n+1} = x(x+y)^n + y(x+y)^n
            \]
            shows that \((x+y)^{n+1}\) is also a polynomial in \(x\) and \(y\), and
            \[
                [(x+y)^{n+1}]_{j,k}
                =
                [(x+y)^n]_{j-1,k}
                +
                [(x+y)^n]_{j,k-1}.
            \]

            If \(j+k=n+1\), then \((j-1)+k=n\) and \(j+(k-1)=n\). Hence the right-hand side is
            \[
                \binom{n}{k} + \binom{n}{k-1}
                =
                \binom{n+1}{k},
            \]
            by Pascal’s identity.

            If \(j+k \ne n+1\), then \((j-1)+k \ne n\) and \(j+(k-1) \ne n\), so both terms are zero.

            Thus,
            \[
                (x+y)^{n+1}
                =
                \sum_{k=0}^{n+1}
                \binom{n+1}{k}
                x^{\,n+1-k} y^{\,k},
            \]
            which is the inductive hypothesis with \(n+1\) substituted for \(n\), completing the
            inductive step.
        \end{proof}

        % Problem 16
        \textbf{16.} Prove Vandermonde's Identity:
        \begin{displaymath}
            \binom{m+n}{r} = \sum_{k=0}^r \binom{m}{k}\binom{n}{r-k}
        \end{displaymath}
        \begin{proof}
            The left-hand side counts all the $k$-element subsets of $[m + n]$.
            The right-hand side does the same thing according to the number of elements chosen from $[n]$.
            Indeed, we can first choose $k$ elements from $[m]$ in $\binom{m}{k}$ ways, then choose the remaining
            $r - k$ elements from the set $\{m+1, m+2, \ldots, m+n \}$ in $\binom{n}{r-k}$ ways.
        \end{proof}

        % Problem 19
        \textbf{19.} Prove the identity:
        \begin{displaymath}
            \binom{n+1}{r+1} = \sum_{k=r}^n \binom{k}{r}
        \end{displaymath}
        \begin{proof}
            Let's go ahead and expand out the right hand side here so we see that:
            \begin{displaymath}
                \binom{n+1}{r+1} = \binom{r}{r} + \binom{r+1}{r} + \binom{r+2}{r} + \cdots + \binom{n}{r}
            \end{displaymath}
            Now with regards to the identity in question, the left-hand side clearly counts all the $r+1$-element subsets
            of $[n+1]$.
            And if we separate the right-hand side into cases according to the largest element.
            That is, there are $\binom{r}{r}$ subsets of $[n+1]$ that have $r + 1$ elements whos largest element is
            $r + 1$; there are $\binom{r+1}{r}$ subsets of $[n+1]$ that have $r+1$ elements whos largest element is
            $r +2$ and so on.
            In general there are $\binom{r+k}{r}$ subsets of $[n+1]$ that have $r+1$ elements whos largest element
            is $r+k+1$ for all $k\le n -r$.
        \end{proof}

        % Problem 20
        \textbf{20.} Prove the identity:
        \begin{displaymath}
            2^n - 1 = \sum_{k=0}^{n-1} 2^k
        \end{displaymath}
        \begin{proof}
            We begin by noticing that the left hand side counts the non-empty subsets of $[n]$.
            To see the count of the right-hand side we count by classifying all the subsets of $A\subseteq[n]$ by each's
            largest element $m\in A$ where $m\in\{1,2,\ldots,n\}$, $\text{max}(A) = m$.
            We notice that every element of $A$ is chosen from $[n]$, and elements $\{1, 2, \ldots, m-1\}$  can
            be included or not independently. So we have that:
            \begin{displaymath}
                2^n - 1 = \sum_{m=1}^{n} 2^{m-1}.
            \end{displaymath}
            A simple shift in indices $k=m-1$ then results in:
            \begin{displaymath}
                2^n - 1 = \sum_{k=0}^{n-1} 2^k.
            \end{displaymath}

        \end{proof}

        % Problem 22
        \textbf{22.} Prove the identity:
        \begin{displaymath}
            \binom{n}{k}\binom{k}{t} = \binom{n}{t}\binom{n-t}{k-t}
        \end{displaymath}
        \begin{proof}
            The left-hand side is the number of ways that we can choose $S$ of $k$-elements from $[n]$, and then
            subsequetly choose $t$-elements from $S$.
            Regarding counting the right-hand side we see that we are immediately choosing $t$-elements from our
            original set of $[n]$.
            But then we are left with a set of $[n-t]$ from which we need to choose $k-t$ elements.
            So it works.
        \end{proof}

        % Problem 24
        \textbf{24.} Prove the identity:
        \begin{displaymath}
            \sum_{k=0}^m\binom{m}{k}\binom{n}{r+k} = \binom{m+n}{m+r}
        \end{displaymath}
        \begin{proof}
            Let $A$ and $B$ be disjoint sets with $|A|=m$ and $|B|=n$.
            We want to find the number of ways to chose $S\subseteq A\cup B$ with $|S| = m+r$, clearly
            $\binom{m+n}{m+r}$.

            Now to find the right-hand side we count by spltting according to how many are chosen from $A$.
            For any such subset $S$, let $k=|S\cap A|$. We see the following
            \begin{itemize}
                \item $j\in\{0,\ldots,m\}$, and
                \item once $j$ is fixed we must choose:
                \begin{itemize}
                    \item $j$-elements from $A$: $\binom{m}{j}$ ways, and
                    \item the remaining $((m+r)-j)$-elements from $B$ in $\binom{n}{(m+r)-j}$ ways.
                \end{itemize}
            \end{itemize}
            So the total count follows as:
            \begin{displaymath}
                \sum_{j=0}^m\binom{m}{j}\binom{n}{(m+n)-j}
            \end{displaymath}
            Now we re-index setting $k=m-j$ yielding $(m+n)-j = r+k$ with $\binom{m}{j}=\binom{m}{m-j} = \binom{m}{k}$,
            and sum becomes:
            \begin{displaymath}
                \sum_{k=0}^m\binom{m}{k}\binom{n}{m+k}
            \end{displaymath}
            and we are done.
        \end{proof}

        % Problem 26
        \textbf{26.} Prove the identity:
        \begin{displaymath}
            \sum_{k=0}^n k^2\binom{n}{k} = n(n-1)2^{n-2} + n2^{n-1} = n(n+1)2^{n-2}
        \end{displaymath}
        \begin{proof}
            Begin with the binomial theorem for $(x+1)^n$, which follows as:
            \begin{displaymath}
                (x+1)^n = \sum_{k=0}^n \binom{n}{k}x^k
            \end{displaymath}
            Now differentiate both sides, and we get
            \begin{displaymath}
                n(x+1)^{n-1} = \sum_{k=1}^n k\binom{n}{k}x^{k-1}
            \end{displaymath}
            Now multiply both sides of the equation by $x$ yielding:
            \begin{displaymath}
                nx(1+x)^{n-1} = \sum_{k=1}^n k\binom{n}{k}x^{k}
            \end{displaymath}
            Differentiating again we get
            \begin{displaymath}
                n[(x+1)^{n-1} + (n-1)x(1+x)^{n-2}] = \sum_{k=1}^n k^2\binom{n}{k}x^{k-1}
            \end{displaymath}
            Now if we substitute in $x=1$ we get
            \begin{eqnarray*}
                n[2^{n-1} + (n-1)x2^{n-2}] & = & \sum_{k=1}^n k^2\binom{n}{k} \\
                n(n+1)2^{n-2} & = & \sum_{k=1}^n k^2\binom{n}{k}
            \end{eqnarray*}
        \end{proof}

        % Problem 33
        \textbf{33.} Prove the identity:
        \begin{displaymath}
            \sum_{k=n-s}^{m-r} \binom{m-k}{r}\binom{n+k}{s} = \binom{m+n+1}{r+s+1}
        \end{displaymath}
        \begin{solution}
        \end{solution}
    \end{onehalfspace}
\end{document}
